\documentclass[11pt,a4paper]{article}

\def\nyear {2026}
\def\nterm {Spring}
\def\nlecturer {Chen Meiri}
\def\ncourse {Noncommutative algebra}

\makeatletter

% packages
\usepackage{amssymb,amsfonts,amsmath,calc,tikz,pgfplots,geometry,mathtools}
\usepackage{color}   % May be necessary if you want to color links
\usepackage[svgnames]{xcolor}
\usepackage[hidelinks]{hyperref}
\usepackage{forest}
\usepackage{commath} % ew
\usepackage{esvect} % \vv makes vectors
\usepackage{centernot} % for the comparison
\usepackage{esdiff}
\usepackage{esint}
\usepackage{amsthm}
\usepackage{fancyhdr}
\usepackage{bm}
\usepackage{witharrows}
\usepackage{bookmark}
\usepackage{tikz-cd}
\usepackage{quiver}
\usepackage{bbm}
\usepackage{textcomp}
\usepackage{float}
\usepackage{gensymb}
\usepackage{cleveref}
\usepackage{environ}
\usepackage{comment}
\usepackage{cancel}
\usepackage{xparse}

% Inevitable cs
\usepackage{listings}
\usepackage[]{algorithm2e}

% tikz libraries
\usetikzlibrary{positioning}
\usetikzlibrary{matrix}
\usetikzlibrary{arrows}
\usetikzlibrary{bending}
\usetikzlibrary{arrows.meta}
\usetikzlibrary{decorations.markings}
\usetikzlibrary{decorations.pathreplacing}
\usetikzlibrary{decorations.markings}
\usetikzlibrary{patterns}
\usetikzlibrary{backgrounds}

% Page style setup
\pagestyle{fancy}
\fancyhfoffset{0pt}
\renewcommand{\headwidth}{\textwidth}
\geometry{margin=1in}
\pgfplotsset{compat=1.18}
\setlength{\headheight}{14.6pt}
\addtolength{\topmargin}{-1.6pt}
\hypersetup{
    colorlinks=false,
    linktoc=section,
    linkcolor=black,
}

%% maketitle setup
\ifx \nauthor\undefined
  \def\nauthor{Yeheli Fomberg}
\else
\fi

\ifx \nemail\undefined
  \def\nemail{\href{mailto:yp@campus.technion.ac.il}
  {\texttt{<yp@campus.technion.ac.il>}}}
\else
\fi

\ifx \ncoursehead \undefined
  \def\ncoursehead{\ncourse}
\fi

\lhead{\emph{\nouppercase{\leftmark}}}
\ifx \nextra \undefined
  \rhead{
    \ifnum\thepage=1
    \else
      \ncoursehead
    \fi}
\else
  \rhead{
    \ifnum\thepage=1
    \else
      \ncoursehead \ (\nextra)
    \fi}
\fi

\def\notes{notes}
\def\problems{problems}

\ifx \ntype \undefined
  \def\ntype{notes}
\else
\fi

\ifx \ntype \notes
  \let\@real@maketitle\maketitle
  \renewcommand{\maketitle}{\@real@maketitle\begin{center}
  \begin{minipage}[c]{0.9\textwidth}\centering\footnotesize
  These notes are not endorsed by the lecturers.
  I have revised them outside lectures to incorporate supplementary
  explanations, clarifications, and material for fun.
  While I have strived for accuracy, any errors or misinterpretations 
  are most likely mine.
  \end{minipage}\end{center}}
\else
  \ifx \ntype \problems
    \let\@real@maketitle\maketitle
    \renewcommand{\maketitle}{\@real@maketitle\begin{center}
    \begin{minipage}[c]{0.9\textwidth}\centering\footnotesize
    These problems are mostly based on the problem sets I was given
    when taking the course.
    
    While I have strived for accuracy and completeness,
    some of the problems would not have been here without the help of
    friends.
    \end{minipage}\end{center}}
  \else
  \fi
\fi


% theorem environments
\theoremstyle{definition}
\newtheorem{definition}{Definition}[section]
\newtheorem{remark}{Remark}[section]
\newtheorem{example}{Example}[section]
\newtheorem{exercise}{Exercise}[section]
\newtheorem{paradox}{Paradox}[section]
\newtheorem{entry}{Entry}[section]
\newtheorem*{solution}{Solution}
\newtheorem*{question}{Question}
\newtheorem*{answer}{Answer}
\newtheorem*{problem}{Problem}
\theoremstyle{plain}
\newtheorem{theorem}{Theorem}[section]
\newtheorem{proposition}[theorem]{Proposition}
\newtheorem{lemma}[theorem]{Lemma}
\newtheorem{corollary}[theorem]{Corollary}
\newenvironment{proofof}[1]{%
    \renewcommand{\proofname}{Proof of #1}%
    \begin{proof}
}{%
    \end{proof}
}
% tikz customization
\pgfarrowsdeclarecombine{twolatex'}{twolatex'}{latex'}{latex'}{latex'}{latex'}
\tikzset{->/.style = {decoration={markings,
                                  mark=at position 0.999
                                  with {\arrow[scale=2]{latex'}}},
                      postaction={decorate}}}
\tikzset{<-/.style = {decoration={markings,
                                  mark=at position 0.001 with {\arrowreversed[scale=2]{latex'}}},
                      postaction={decorate}}}
\tikzset{-->/.style = {decoration={markings,
                                  mark=at position 0.999
                                  with {\arrow[scale=2]{latex'[sep=0pt -2.5]}}},
                      postaction={decorate}}}
\tikzset{<--/.style = {decoration={markings,
                                  mark=at position 0.001
                                  with {\arrowreversed[scale=2]{latex'[sep=0pt -2.5]}}},
                      postaction={decorate}}}
\tikzset{<->/.style = {decoration={markings,
                                  mark=at position 0.001 with {\arrowreversed[scale=2]{latex'}},
                                  mark=at position 0.999 with {\arrow[scale=2]{latex'}},},
                      postaction={decorate}}}
\tikzset{->-/.style = {decoration={markings,
    mark=at position 0.50+0.225/#1 with {\arrow[scale=2]{latex'}}},
                      postaction={decorate}}}
\tikzset{-<-/.style = {decoration={markings,
                                   mark=at position 0.50-0.225/#1 with {\arrowreversed[scale=2]{latex'}}},
                       postaction={decorate}}}
\tikzset{->>/.style = {decoration={markings,
                                  mark=at position 1 with {\arrow[scale=2]{latex'}}},
                      postaction={decorate}}}
\tikzset{<<-/.style = {decoration={markings,
                                  mark=at position 0 with {\arrowreversed[scale=2]{twolatex'}}},
                      postaction={decorate}}}
\tikzset{<<->>/.style = {decoration={markings,
                                   mark=at position 0 with {\arrowreversed[scale=2]{twolatex'}},
                                   mark=at position 1 with {\arrow[scale=2]{twolatex'}}},
                       postaction={decorate}}}
\tikzset{->>-/.style = {decoration={markings,
                                   mark=at position 0.50+0.450/#1 with {\arrow[scale=2]{twolatex'}}},
                       postaction={decorate}}}
\tikzset{-<<-/.style = {decoration={markings,
                                   mark=at position 0.50-0.450/#1 with {\arrowreversed[scale=2]{twolatex'}}},
                       postaction={decorate}}}

\pgfarrowsdeclare{biggertip}{biggertip}{%
  \setlength{\arrowsize}{1pt}
  \addtolength{\arrowsize}{.1\pgflinewidth}
  \pgfarrowsrightextend{0}
  \pgfarrowsleftextend{-5\arrowsize}
}{%
  \setlength{\arrowsize}{1pt}
  \addtolength{\arrowsize}{.1\pgflinewidth}
  \pgfpathmoveto{\pgfpoint{-5\arrowsize}{4\arrowsize}}
  \pgfpathlineto{\pgfpointorigin}
  \pgfpathlineto{\pgfpoint{-5\arrowsize}{-4\arrowsize}}
  \pgfusepathqstroke
}
\tikzset{
	EdgeStyle/.style = {>=biggertip}
}

\tikzset{circ/.style = {fill, circle, inner sep = 0, minimum size = 3}}
\tikzset{scirc/.style = {fill, circle, inner sep = 0, minimum size = 1.5}}
\tikzset{mstate/.style={circle, draw, black, text=black, minimum width=0.7cm}}

\tikzset{eqpic/.style={baseline={([yshift=-.5ex]current bounding box.center)}}}

\definecolor{mblue}{rgb}{0.2, 0.3, 0.8}
\definecolor{morange}{rgb}{1, 0.5, 0}
\definecolor{mgreen}{rgb}{0, 0.4, 0.2}
\definecolor{mred}{rgb}{0.5, 0, 0}

% topology
\DeclareMathOperator{\dfr}{dfr} % deformation retract
\newcommand{\Cells}{\text{Cells}}
\newcommand{\RP}{\mathbb{RP}}
\newcommand{\CP}{\mathbb{CP}}

% order theory
\newcommand{\join}{\vee}
\newcommand{\meet}{\wedge}

% algebraic geometry
\DeclareMathOperator{\Rad}{Rad} % Radical of modules
\DeclareMathOperator{\Spec}{Spec}
% \renewcommand{\div}{\mathrm{div}} % divisors
\DeclareMathOperator{\Mer}{Mer} % Meromorphic functions
\DeclareMathOperator{\Exc}{Exc} % Exceptional divisor
\DeclareMathOperator{\Td}{Td} % Todd classes

% algebra

\DeclareMathOperator{\trdeg}{tr.deg.}
\DeclareMathOperator{\odd}{odd}
\DeclareMathOperator{\even}{even}
\DeclareMathOperator{\lcm}{lcm}
\DeclareMathOperator{\ord}{ord}
\DeclareMathOperator{\codim}{codim}
\DeclareMathOperator{\Ann}{Ann}
\DeclareMathOperator{\Ext}{Ext}
\DeclareMathOperator{\Sym}{Sym}
\DeclareMathOperator{\Out}{Out}
\DeclareMathOperator{\Aut}{Aut}
\DeclareMathOperator{\End}{End}
\DeclareMathOperator{\Inn}{Inn}
\DeclareMathOperator{\Mat}{Mat}
\DeclareMathOperator{\Fix}{Fix}
\DeclareMathOperator{\std}{std}
\DeclareMathOperator{\sgn}{sgn}
\DeclareMathOperator{\adj}{adj}
\DeclareMathOperator{\id}{id}
\DeclareMathOperator{\op}{op}
\DeclareMathOperator{\tr}{tr}
\DeclareMathOperator{\diag}{diag}
\DeclareMathOperator{\rank}{rank}
\DeclareMathOperator{\rk}{rk}
\DeclareMathOperator{\coker}{coker}
\DeclareMathOperator{\Mult}{Mult} % multiplier algebra
\DeclareMathOperator{\Frac}{Frac} % fraction field
\DeclareMathOperator{\Rat}{Rat}
\DeclareMathOperator{\Gal}{Gal}
\newcommand{\ioplus}{\overset{\text{int}}{\oplus}} % interior direct sum
\newcommand{\GG}{\mathcal G} % Galois correspondence
\newcommand{\FF}{\mathcal F} % Galois correspondence
\newcommand{\cupp}{\smile} % cup product
\newcommand{\capp}{\frown} % cap product
\newcommand{\actson}{\curvearrowright}
\newcommand{\bra}{\left\langle}
\newcommand{\ket}{\right\rangle}
\newcommand{\idealin}{\triangleleft\,}
\newcommand{\normalin}{\triangleleft\,}
\newcommand{\ip}[2]{\left\langle #1, #2 \right\rangle}
\newcommand{\bigslant}[2]
{{\raisebox{.2em}{$#1$}\left/\raisebox{-.2em}{$#2$}\right.}}
\newcommand{\properideal}{%
  \mathrel{\ooalign{$\lneq$\cr\raise.22ex\hbox{$\lhd$}\cr}}}

% categories
\newcommand{\cat}[1]{\mathbf{#1}}
\newcommand{\Ch}{\cat{Ch}} % Chain complexes   
\newcommand{\Set}{\cat{Set}}
\newcommand{\Grp}{\cat{Grp}}
\newcommand{\Ab}{\cat{Ab}}
\newcommand{\Ring}{\cat{Ring}}
\newcommand{\CRing}{\cat{CRing}}
\newcommand{\Top}{\cat{Top}}
\newcommand{\Vect}{\cat{Vect}}
\newcommand{\cmod}{\cat{mod}}

%% matrix groups
\newcommand{\GL}{\mathrm{GL}}
\newcommand{\Or}{\mathrm{O}}
\newcommand{\PGL}{\mathrm{PGL}}
\newcommand{\PSL}{\mathrm{PSL}}
\newcommand{\PSO}{\mathrm{PSO}}
\newcommand{\PSU}{\mathrm{PSU}}
\newcommand{\SL}{\mathrm{SL}}
\newcommand{\msl}{\mathrm{sl}}
\newcommand{\SO}{\mathrm{SO}}
\newcommand{\Spin}{\mathrm{Spin}}
\newcommand{\Sp}{\mathrm{Sp}}
\newcommand{\SU}{\mathrm{SU}}
\newcommand{\U}{\mathrm{U}}
\let\char\relax
\newcommand{\char}{\text{char}}

% analysis
\renewcommand{\Re}{\mathrm{Re}}
\renewcommand{\Im}{\mathrm{Im}}
\NewDocumentCommand{\dx}{o}{
  \IfNoValueTF{#1}
    {\dif x}                  % if no optional argument
    {\dif^{#1} \!x}         % if optional argument
}
\NewDocumentCommand{\dy}{o}{
  \IfNoValueTF{#1}
    {\dif y}                  % if no optional argument
    {\dif^{#1} \!y}         % if optional argument
}
\newcommand{\dt}{\dif t}
\newcommand{\du}{\dif u}
\newcommand{\dv}{\dif v}
\newcommand{\dz}{\dif z}
\newcommand{\ds}{\dif s}
\newcommand{\dw}{\dif w}
\newcommand{\dr}{\dif r}
\newcommand{\dm}{\dif m}
\newcommand{\df}{\dif f}
\newcommand{\dV}{\dif V}
\newcommand{\dS}{\dif S}
\newcommand{\dA}{\dif \mathrm{A}}
\newcommand{\dmu}{\dif \mu}
\newcommand{\dnu}{\dif \nu}
\newcommand{\dtheta}{\dif \theta}
\newcommand{\domega}{\dif \omega}
\newcommand{\dsigma}{\dif \sigma}
\newcommand{\dtau}{\dif \tau}
\newcommand{\dvarphi}{\dif \varphi}
\renewcommand{\dif}{\mathrm{d}}
\renewcommand{\Dif}{\mathrm{D}}
\renewcommand{\triangle}{\Delta}
\newcommand{\ptriangle}{\partial \triangle}
\DeclareMathOperator{\im}{im}
\DeclareMathOperator{\cis}{cis}
\DeclareMathOperator{\rot}{rot}
\DeclareMathOperator{\vspan}{span}
\DeclareMathOperator{\grad}{grad}
\DeclareMathOperator{\curl}{curl}
\DeclareMathOperator{\esssup}{esssup}
\newcommand{\hodge}{{\mathnormal{\star}}}
\DeclareMathOperator{\range}{range}
\DeclareMathOperator{\Hom}{Hom}
\DeclareMathOperator{\Int}{Int}
\DeclareMathOperator{\Conv}{Conv} % convex hull
\newcommand{\ind}{\mathbbm{1}}
\DeclareMathOperator{\Res}{Res}
\DeclareMathOperator{\graph}{graph}
\DeclareMathOperator{\diam}{diam}
\DeclareMathOperator{\supp}{supp}
\DeclareMathOperator{\Vol}{Vol} % Volume
\DeclareMathOperator{\Area}{Area} % Area
\DeclareMathOperator{\area}{area}
\DeclareMathOperator{\Ind}{Ind} % Index
\newcommand{\length}{\mathrm{length}}

% logic
\DeclareMathOperator{\MOD}{MOD}
\DeclareMathOperator{\Theory}{Theory}
\newcommand{\notimplies}{%
  \mathrel{{\ooalign{\hidewidth$\not\phantom{=}$\hidewidth\cr$\implies$}}}}

% nice
\newcommand{\half}{\frac{1}{2}}
\newcommand{\pair}{\del}
\newcommand{\taking}[1]{\xrightarrow{#1}}
\newcommand{\otaking}[1]{\xleftarrow{#1}}
\newcommand{\inv}{^{-1}}
\newcommand{\ot}{\leftarrow}
\newcommand{\ninfty}{-\infty}
\newcommand{\pinfty}{+\infty}
\newcommand{\floor}[1]{\left\lfloor #1 \right\rfloor}
\newcommand{\ceil}[1]{\left\lceil #1 \right\rceil}
\newcommand{\viff}{\Big\Updownarrow} % vertical iff
\newcommand{\bmid}{\,\middle\vert\,} % big mid

% probability
\newcommand{\Prob}{\mathbf{P}}
\renewcommand{\vec}[1]{\boldsymbol{\mathbf{#1}}}
\DeclareMathOperator{\Bin}{Bin}
\DeclareMathOperator{\Geo}{Geo}
\DeclareMathOperator{\Poi}{Poi}
\DeclareMathOperator{\Exp}{Exp}
\DeclareMathOperator{\Var}{Var} % Variance
\DeclareMathOperator{\Cov}{Cov}

% special letters
\newcommand{\N}{\mathbb{N}}
\newcommand{\Z}{\mathbb{Z}}
\newcommand{\Q}{\mathbb{Q}}
\newcommand{\R}{\mathbb{R}}
\newcommand{\C}{\mathbb{C}}
\newcommand{\F}{\mathbb{F}}
\newcommand{\E}{\mathbf{E}}
\newcommand{\D}{\mathbb{D}}
\renewcommand{\H}{\mathbb{H}}
\newcommand{\ps}{\mathcal{P}}
\newcommand{\M}{\mathcal{M}}
\renewcommand{\L}{\mathcal{L}}
\renewcommand{\O}{\mathcal{O}}
\renewcommand{\U}{\mathcal{U}}
\newcommand{\Omicron}{O}
\newcommand{\powerset}{\mathcal{P}}

% text
\newcommand{\st}{\text{ s.t. }}
\newcommand{\tand}{\quad \text{and} \quad}
\newcommand{\tor}{\quad \text{or} \quad}
\newcommand{\stand}{\text{ and }}
\newcommand{\stor}{\text{ or }}
\renewcommand{\tt}[1]{\textnormal{\textbf{(#1).}}} %tt=theorem title GET RID OF

% title format
\title{\textbf{\ncourse}}

\ifx \ntype \notes
  \author{Notes taken by \nauthor\, \nemail \\ Based on lectures by \nlecturer}
  \date{\nterm\ \nyear}
\else
  \ifx \ntype \problems
    \author{These problems were solved by \nauthor \\ \nemail}
  \else
  \fi
\fi

\makeatother


\begin{document}
\maketitle

% Insert cool image here

\newpage
\tableofcontents
\newpage

\section{Introduction}

\begin{remark}
  In this course, unless explicitly stated otherwise, the term
  ring will be preserved to rings with units which are not necessarily
  commutative.
  Homomorphisms between rings are required to respect the identity,
  that means that if $f \colon R \to S$ is a ring homomorphism
  then $f(1_R) = 1_S$.
\end{remark}

\begin{definition}[Opposite ring]
  For a ring $R$, we define the opposite ring $R^{\op}$ by giving
  it the additive ring of $R$, and defining $a *_{R^{\op}} b = b *_R a$.
\end{definition}

\begin{example}
  Let $n \geq 1$, and let $\mathbb F$ be a field.
  The map $A \mapsto A^T$ is an isomorphism from $M_n(\mathbb F)$ to $M_n(\mathbb F)^{\op}$.
\end{example}

\subsection{Modules}

\begin{definition}[Module]
  Let $R$ be a ring.
  A left $R$-module is a pair $(M,\rho)$ where $M$ is an abelian group,
  and $\rho$ is a homomorphism from $R$ to $\End(M)$.
  Similarly, a right $R$-module is a pair $(M,\rho)$ where $M$ is an abelian
  group, and $\rho$ is a homomorphism from $R$ to $\End(M)^{\op}$.
\end{definition}

\begin{remark}
  If $R$ is commutative, then every left $R$-module is also a right $R$-module
  and vice versa.
  Thus, when we do commutative algebra, we simply use the term ``module''.
\end{remark}

\begin{remark}
  We sometimes say that $M$ is a left (or right) module,
  without mentioning the homomorphism $\rho$ when it is possible to infer
  it from context.
  If $r \in R$ and $m \in M$ we usually write $r \cdot m$ instead
  of $\phi(r)(m)$.
  This terminology is justified by the following exercise.
\end{remark}

\begin{exercise}
  Let $R$ be a ring and $M$ and abelian group.
  Prove that there is a bijection between the set of left $R$-modules
  with underlying group $M$ and the set of
  functions $\cdot \colon R \times M \to M$ which satisfy the following
  properties:
  \begin{enumerate}
    \item[(1)] $1_R \cdot m = m$.
    \item[(2)] $r \cdot (m + n) = r \cdot m + r \cdot n$.
    \item[(3)] $(r + s) \cdot m = r \cdot m + s \cdot m$.
    \item[(4)] $(rs) \cdot m = r \cdot (s \cdot m)$.
  \end{enumerate}
\end{exercise}

\begin{example}
  A vector space is a module over a field.
\end{example}

\begin{example}
  Every abelian group is a module over $\Z$
  (note that there is a unique homomorphism from $\Z$ to $\End(M)$).
\end{example}

\begin{example}
  Let $n,m \geq 2$ and $R$ be a ring.
  Then $M_{n,m}(R)$ is a left $M_n(R)$-module
  where the action of $R$ on $M$ is by matrix multiplication from the left.
  Similarly $M_{m,n}(R)$ is a right $M_n(R)$-module where the action of $R$
  on $M$ is by matrix multiplication from the right.
\end{example}

\begin{example}
  Let $R$ be a ring.
  Then $R$ is a left (equivalently, right) $R$-module where the action of $R$
  on $R$ is by left (equivalently, right) multiplication.
  We usually write $_R R$ (equivalently $R_R$) to indicate that we view
  $_R R$ as a left module (equivalently, right).
  The left (equivalently, right) ideals of $R$ are the left (equivalently, right) submodules
  of $_R R$ (equivalently, $R_R$).
\end{example}

\begin{remark}
  Many definitions, constructions, and theorems about abelian groups have
  natural counterparts in the theory of modules.
\end{remark}

\begin{definition}
  Let $R$ be a ring.
  \begin{enumerate}
    \item[(1)] A $R$-homomorphism from a left $R$-module $M$
      to a left $R$-module $N$ is a group homomorphism $\alpha \colon M \to N$
      such that for every $r \in R$ and every $m \in M$
      we have that $\alpha(rm) = r \alpha(m)$.
    \item[(2)] A submodule of a left $R$-module $M$ is a subgroup
      $L \leqslant M$ such that for every $r \in R$ and every $m \in L$
      we have that $rm \in L$.
    \item[(3)] Let $L$ be a submodule of a left $R$-module $M$.
      The quotient group $M/L$ inherits a structure of a left $R$-module
      by the rule ``for every $r \in R$ and every $m + L \in M/L$
      we define $r(m + L) := rm + L$''.
  \end{enumerate}
\end{definition}

\begin{example} \phantom{}
  \begin{enumerate}
    \item[(1)] Let $M$ be an abelian group and view it as a $\Z$-module.
      Every subgroup of $M$ is a submodule.
    \item[(2)] Let $R$ be a ring. The submodules of $R_R$ are the right ideals
      of $R$.
  \end{enumerate}
\end{example}

\begin{exercise}
  Let $R$ be a ring and $\alpha \colon M \to N$ be a homomorphism of left (right)
  $R$-modules.
  Prove that $\ker \alpha$ and $\im \alpha$ are left (right) submodules of $M$
  and $N$ respectively.
\end{exercise}

\begin{exercise}
  State and prove the isomorphism theorems for modules.
\end{exercise}

\begin{definition}[Direct product]
  Let $R$ be a ring.
  Let $(M_i)_{i \in I}$ be an indexed family of left $R$-modules.
  The direct product of $(M_i)_{i \in I}$, denoted by $\prod_{i \in I} M_i$,
  is the left $R$-module whose underlying group is the direct product
  of the underlying groups of each of the $M_i$s,
  and the action of $R$ is the pointwise actions.
  More explicitly:
  \begin{enumerate}
    \item[(1)] The elements of $\prod_{i \in I} M_i$
      are functions $f \colon I \to \cup_{i \in I} M_i$ such that
      for every $i \in I$ we have $f(i) \in M_i$.
    \item[(2)] For every $f,g \in \prod_{i \in I} M_i$, the sum $f + g$
      is defined by the pointwise sum $(f + g)(i) := f(i) + g(i)$.
    \item[(3)] For every $r \in R$ and $f \in \prod_{i \in I} M_i$,
      the product $r \cdot f$ is defined by $(r \cdot f)(i) = r \cdot f(i)$
      for evert $i \in I$.
  \end{enumerate}
\end{definition}

\begin{definition}[Direct sum]
  Let $(M_i)_{i \in I}$ be an indexed family of left $R$-modules.
  The direct sum of the $M_i$s denoted by $\oplus_{i \in I} M_i$,
  is the submodule of $\prod_{i \in I} M_i$ consisting of the elements
  $f$ with finite support, i.e., the elements $f \in \prod_{i \in I} M_i$
  for which the set $\{i \in I \mid f(i) \neq 0\}$ is finite.
\end{definition}

\begin{remark}
  We remind that all of the above (and below) 
  definitions, theorems, and exercises
  have counterparts with right $R$-modules.
\end{remark}

\begin{exercise}
  Let $R$ be a ring and $M$ a left or right $R$-module.
  Prove that an intersection of (possibly infinitely many) submodules of $M$
  is a submodule of $M$.
\end{exercise}

\begin{definition}[Generated submodule]
  Let $X$ be a subset of $M$.
  The submodule of $M$ generated by $X$ is the intersection of all submodules
  of $M$ that contain $X$.
\end{definition}

\begin{definition}[Sum of submodules]
  Let $\{N_i \mid i \in I\}$ be a set consisting of submodules of $M$.
  We denote by $\sum_{i \in I} N_i$ the submodule of $M$
  generated by $\cup_{i \in I} N_i$.
\end{definition}

\begin{exercise}
  Let $R$ be a ring, and let $M$ be a left or right $R$-module.
  Prove that the submodule generated by $X$ consists of the $R$-linear
  finite sums of elements of $X$ (where the empty sum is defined to be
  equal to zero).
\end{exercise}

\begin{definition}
  Let $R$ be a ring, and let $M$ be a left or right $R$-module.
  We define $\End_R(M)$ to be the subring of $\End(M)$ consisting of the
  $R$-homomorphisms.
\end{definition}

\begin{exercise}
  Let $R$ be a commutative ring and $M$ an $R$-module.
  Prove that for every $r \in R$ and every $f \in \End_R(M)$,
  the function $rf \colon M \to M$ defined by $(rf)(m) := r(f(m))$
  belongs to $\End_R(M)$.
  Deduce that in this case, $\End_R(M)$ carries the structure of an $R$-module.
\end{exercise}

\begin{definition}[Noetherian ring]
  Let $R$ be a ring, and let $M$ be a left or right module.
  The module $M$ is called Noetherian (respectively Artinian)
  if there is no infinite and strictly increasing (respectively, decreasing)
  chain of submodules of $M$.
\end{definition}

\begin{example}
  $\Z_\Z$ is a Noetherian $\Z$-module which is not Artinian.
\end{example}

\begin{example}
  \label{ex:prufer}
  Let $p$ be a prime number.
  Then $\Z[1/p]/\Z = \{a/p^n \mid a \in \Z \stand n \in \N\}/\Z$
  is an Artinian $\Z$-module which is not Noetherian.
\end{example}

\begin{remark}
  The group $\Z[1/p]/\Z$ is known as a Pr\"ufer group,
  and is usually denoted by $\Z(p^{\infty})$.
  It is worth taking the time to understand its structure,
  and why the claim in \Cref{ex:prufer} holds.
\end{remark}

\subsection{Simple and semisimple modules}

\begin{definition}[Simple module]
  Let $R$ be a ring.
  A left or right $R$-module $M$ is called simple if $M$ is not the zero
  module and the only submodules of $M$ are $0$ and $M$.
\end{definition}

\begin{example}
  A finite abelian group is a simple module over $\Z$ if and only if it is
  of a prime order.
\end{example}

\begin{example}
  Let $n,m \geq 1$ and $\mathbb F$ be a field.
  Then $M_{n \times m}(\mathbb F)$ is a simple left $M_n(R)$-module
  if and only if $m = 1$.
\end{example}

\begin{example}
  Let $R$ be a nonzero ring.
  Then $R_R$ is a simple right $R$-module if and only if the right ideals
  of $R$ are the zero ideal and $R$ itself.
\end{example}

\begin{exercise}[Characterization of simple modules]
  Let $R$ be a ring.
  For a nonzero right $R$-module $M$,
  show that the following conditions are equivalent:
  \begin{enumerate}
    \item[(1)] $M$ is simple.
    \item[(2)] $mR = M$ for all nonzero $m \in M$.
    \item[(3)] $M \cong R_R/I$ for some maximal right ideal $I$ of $R$.
  \end{enumerate}
\end{exercise}

\begin{lemma}[Schur's lemma]
  Let $M$ and $N$ be right $R$-modules,
  and let $\rho \colon M \to N$ be a nonzero homomorphism.
  Then:
  \begin{enumerate}
    \item[(1)] If $M$ is simple, then $\rho$ is injective.
    \item[(2)] If $N$ is simple, then $\rho$ is surjective.
    \item[(3)] If $M$ and $N$ are simple right $R$-modules,
      then either $M \cong N$ or $\Hom_R(M,N) = 0$.
    \item[(4)] If $M$ is simple, then $\End_R(M)$ is a division ring.
  \end{enumerate}
\end{lemma}

\begin{remark}
  Recall that a division ring is a ring in which division is well defined for nonzero
  elements.
  A commutative division ring is a field.
  As a fun fact, Wedderburn's little theorem asserts that all finite division rings
  are commutative and therefore finite fields.
\end{remark}

\begin{definition}[Composition series]
  Let $R$ be a ring and $M$ be a left or right $R$-module.
  An increasing chain of submodules
  $0 = M_0 \subsetneq M_1 \subsetneq \cdots \subsetneq M_{n-1} \subsetneq M_n = M$ 
  of $M$ is called a composition series of $M$ if for every $0 \le i < n$,
  the $R$-module $M_{i+1}/M_i$ is simple.
  The factor modules are called the composition modules of $M$.
\end{definition}

\begin{example}
  The $\Z$-module $\Z_{\Z}$ does not have a composition series.
\end{example}

\begin{example}
  For every natural number $n \geq 2$ denote $\overline \Z_n := \Z/n\Z$.
  The series $0 \subsetneq 4 \overline \Z_{12} \subsetneq 2 \overline \Z_{12} \subsetneq
  \overline \Z_{12}$ and
  $0 \subsetneq 6 \overline \Z_{12} \subsetneq 3 \overline \Z_{12} \subsetneq
  \overline \Z_{12}$ are both composition series for the $\Z$-module $\overline \Z_{12}$.
  The composition factors of the above two series consist of two simple modules isomorphic
  to $\overline \Z_{2}$ and one simple module isomorphic to $\overline \Z_{3}$.
  The fact that the composition factors are identical in both series is not a coincidence.
\end{example}

\begin{exercise}
  Prove that $M$ has a composition series if and only if it is Noetherian
  and also Artinian.
\end{exercise}

\begin{exercise}
  Prove that Jordan--H\"older theorem:
  If $M$ has a composition series then, up to isomorphism,
  the composition factors and their multiplicity are the same for each
  composition series.
\end{exercise}

\begin{definition}[Semisimple module]
  A right or left $R$-module $M$ is called a semisimple module if it is isomorphic
  to a direct sum of simple modules.
\end{definition}

\begin{example}
  Let $p_1,\dots,p_k$ be prime numbers and denote $n := p_1 \cdots p_k$.
  Then $\Z/n\Z$ is a semisimple $\Z$-module if and only if $p_1,\dots,p_k$
  are distinct.
\end{example}

\begin{definition}[Internal direct sum]
  Let $R$ be a ring and $M$ a right or left $R$-module.
  Let $(M_i)_{i \in I}$ be an indexed collection of submodules of $M$.
  Let $\alpha \colon \oplus_{i \in I} M_i \to \sum_{i \in I} M_i$
  be the $R$-homomorphism defined by $\alpha(f) := \sum_{i \in I} f(i)$
  for every $f \in \oplus_{i \in I} M_i$.
  We say that $\sum_{i \in I} M_i$ is an internal direct sum of the indexes
  collection $(M_i)_{i \in I}$ if $\alpha$ is an isomorphism onto $\sum_{i \in I} M_i$.
\end{definition}

\begin{example}
  The Euclidean plane $M = \R^2$ viewed as an $R$-module is an internal direct sum
  of the submodules $M_1 = \{(x,0) \mid x \in \R\}$ and $M_2 = \{(0,y) \mid y \in \R\}$.
  Indeed, the map $\alpha \colon M_1 \oplus M_2 \to \sum_{i=1}^{2} M_i$ 
  given by $\alpha((x,0),(0,y)) = (x,0) + (0,y) = (x,y)$ is an isomorphism.
\end{example}

\begin{example}
  Any polynomial ring $k[x]$ viewed as a module over $K$ is an internal direct sum
  of the submodules $M_{\even} = \{p \in k[x] \mid p \text{ only contains even powers
  of } x\}$ and $M_{\odd} = \{p \in k[x] \mid p \text{ only contains odd powers
  of } x\}$.
  In that case $\alpha$ is clearly an isomorphism since each polynomial can be uniquely
  decomposed into the sum of the element of each of these modules.
\end{example}

\begin{example}
  The matrix algebra $M_n(\R)$ viewed as an $\R$-module is an internal
  direct sum of the submodules $M_{\mathrm{sym}}$ and $M_{\mathrm{skew}}$
  of symmetric and skew-symmetric matrices by the famous decomposition result
  we have seen in linear algebra.
\end{example}

\begin{remark}
  If $\sum_{i \in I} M_i$ is an internal direct sum,
  we may use the notation $\ioplus_{i \in I} M_i = \sum_{i \in I} M_i$,
  or write $\sum_{i \in I} M_i \overset{\text{int}}{\cong} \oplus_{i \in I} M_i$.
\end{remark}

\begin{lemma}
  \label{lem:ioplus-lemma}
  Let $R$ be a ring, let $M$ be a left or right $R$-module.
  Let $N$ and $L$ be submodules of $M$.
  Then:
  \begin{enumerate}
    \item[(1)] $L + N \overset{\text{int}}{\cong} L \oplus N$ if and only if
      $L \cap N = 0$.
    \item[(2)] If $L$ is simple then $L + N \overset{\text{int}}{\cong} L \oplus N$
      if and only if $L$ is not contained in $N$.
  \end{enumerate}
\end{lemma}
\begin{proof} \phantom{}
  \begin{enumerate}
    \item[(1)] Let $\rho \colon L \oplus N \to L + N$ be the $R$-homomorphism defined by
      $\rho((\ell,n)) = \ell + n$.
      The $R$-homomorphism $\rho$ is always onto, so we only have to check when it is
      injective.
      For every $(\ell,n) \in L \oplus N$, we have that $(\ell,n) \in \ker \rho$
      if and only if $\ell + n = 0$ or equivalently $\ell = -n \in L \cap N$.
      It follows that $\ker \rho = 0$ if and only if $L \cap N = 0$.
    \item[(2)] Since $L$ is simple and $L \cap N$ is a submodule of $L$,
      by definition of the simple module either $L \cap N = 0$
      or $L \cap N = L$.
      By $(1)$, the former holds if and only if
      $L + N \overset{\text{int}}{\cong} L \oplus N$.
      The latter possibility holds if and only if $L \subseteq N$.
      This completes the proof.
  \end{enumerate}
\end{proof}

\begin{remark}
  Part $(1)$ of \Cref{lem:ioplus-lemma} holds even in the case of an infinite
  sums.
\end{remark}

\begin{definition}[Complemented module]
  Let $R$ be a ring and $M$ a left or right $R$-module.
  Then $M$ is called complemented if for every submodule $N$ of $M$
  there exists a submodule $L$ such that $M = N \ioplus L$.
\end{definition}

\begin{proposition}[Properties of complemented modules]
  \label{prop:properties-of-complemented-modules}
  Let $R$ be a ring and $M$ a left or right complemented $R$-module.
  Then:
  \begin{enumerate}
    \item[(1)] Every submodule of $M$ is complemented.
    \item[(2)] If $M \neq 0$ then $M$ contains a simple submodule.
    \item[(3)] $M$ is generated by its simple submodules.
  \end{enumerate}
\end{proposition}
\begin{proof} \phantom{}
  \begin{enumerate}
    \item[(1)] Let $N$ be a submodule of $M$ and $L$ a submodule of $N$.
      We have to show that there exists a submodule $K$ of $N$ such that
      $N = L \ioplus K$.
      Since $M$ is complemented, there exists a submodule $E$ of $M$
      such that $M = N \ioplus E$.
      It follows by \Cref{lem:ioplus-lemma} that $L + E = L \ioplus E$,
      which is a submodule of $M$.
      Since $M$ is complemented, there exists a submodule $D$ of $M$
      such that $M = D \ioplus (L \ioplus E)$.
      Since $L \subseteq N$, we get that:
      \[
        N =
        N \cap M =
        N \cap (D \ioplus (L \ioplus E)) =
        N \cap (L \ioplus (D \ioplus E)) =
        L \ioplus (N \cap (D \ioplus E)),
      \]
      where the second to last transition is commutativity of the direct sum,
      and the last transition is the modular property of the lattice of submodules
      (see \Cref{ssec:lattices}).
      Hence, we can take $K := N \cap (D \ioplus E)$,
      which completes the proof.
    \item[(2)] Assume that $M \neq 0$ and choose some nonzero $m \in M$.
      By Zorn's lemma, there exists a submodule $N$ of $M$ which is maximal
      with respect to the property that it does not contain $m$.
      Since $m \notin N$, we get that $M \neq N$.
      Since $M$ is complemented, there exists a submodule $L$ of $M$
      such that $M = L \ioplus N$.
      We claim that $L$ is simple.
      Assume, by way of contradiction, that $L$ is not simple.
      Since $L$ is complemented (by $(1)$),
      there exist nonzero submodules $L_1$ and $L_2$ of $L$ such that
      $L = L_1 \ioplus L_2$.
      It follows that:
      \[
        (L_1 \ioplus N) \cap
        (L_2 \ioplus N) = N,
      \]
      so $m \notin (L_1 \ioplus N) \cap (L_2 \ioplus N)$.
      Thus, either $m \notin (L_1 \ioplus N)$
      or $m \notin (L_2 \ioplus N)$.
      Since $L_1$ and $L_2$ are nonzero, we get a contradiction to the maximality
      of $N$ with respect to the property of not containing $m$.
    \item[(3)] If $M = 0$ then $M$ is semisimple (this completes the proof).
      Assume that $M \neq 0$.
      Let $N$ be the submodule of $M$ which is generated by all simple
      submodules of $M$.
      Assume, by way of contradiction, that $M \neq N$.
      Since $M$ is complemented, there exists a nonzero submodule $L$
      of $M$ such that $M = L \ioplus N$.
      By $(1)$ we get that $L$ is complemented,
      and by $(2)$ we get that $L$ contains a simple submodule $K$.
      By the definition of $N$, we must have $K \subseteq N \cap L = 0$,
      which is a contradiction (recall a simple submodule is nonzero).
  \end{enumerate}
\end{proof}

\begin{remark}
  Note that the zero module is a semisimple but not simple.
\end{remark}

\begin{theorem}[Characterization of semisimple modules]
  Let $R$ be a ring and $M$ a left or right $R$-module.
  Then the following conditions are equivalent:
  \begin{enumerate}
    \item[(1)] $M$ is semisimple.
    \item[(2)] $M$ is an internal direct sum of some simple submodules of it.
    \item[(3)] $M = \sum_{i \in I} M_i$ where $\{M_i \mid i \in I\}$
      is the set of simple submodules of $M$.
    \item[(4)] $M$ is complemented.
  \end{enumerate}
\end{theorem}
\begin{proof}
  We can assume that $M \neq 0$ (because in that case the statement is trivial).
  The implications $(1) \implies (2) \implies (3)$ are clear.
  We will now show $(3) \implies (4)$.
  Let $N$ be a submodule of $M$.
  Define:
    $T := \{J \subseteq I \mid
      N + \sum_{i \in I} M_i \overset{\text{int}}{\cong}
      N \oplus (\oplus_{i \in I} M_i)\}$.
  $T$ is not empty since $\emptyset \in T$.
  Moreover, $T$ is closed to the union of ascending chains of sets in $T$.
  Thus we can apply Zorn's lemma and get that there exists a maximal
  element $J \in T$ with respect to inclusion.
  It is now enough to show that $M = N \ioplus (\ioplus_{i \in I} M_i)$,
  since them $M$ is complemented which completes the proof.
  Assume otherwise.
  By $(3)$, there exists $j \in I$ such that $N_j$ is not contained in
  $N \ioplus (\ioplus_{i \in I} M_i)$,
  in particular, $j \notin J$.
  \Cref{lem:ioplus-lemma} implies that:
  \[
    N = \sum_{i \in J \cup \set{j}} M_i =
    (N \ioplus (\ioplus_{i \in J \cup \set{j}} M_i)) \ioplus N_j =
    N \ioplus (\ioplus_{i \in J \cup \set{j}} M_i),
  \]
  in contradiction to the maximality of $J$.
  This contradiction completes the proof of the implication $(3) \implies (4)$.
  Note that in the proof of $(3) \implies (4)$, if we choose $N = 0$,
  then we have shown that $M = \ioplus_{i \in I} M_i$,
  which means that $M$ is semisimple hence the implication $(3) \implies (1)$
  is now proved.
  We now have the implications $(1) \iff (2) \iff (3) \implies (4)$,
  but the implication $(4) \implies (3)$ follows from
  \Cref{prop:properties-of-complemented-modules}.
\end{proof}

\begin{corollary}
  A submodule of a semisimple module is semisimple.
\end{corollary}
\begin{proof}
  This follows immediately from the characterization of semisimple modules as
  complemented modules.
  Since any submodule of a complemented module is itself complemented,
  the proof is complete.
\end{proof}

\begin{exercise}
  Let $M$ be a semisimple module.
  Then the following are equivalent:
  \begin{enumerate}
    \item[(1)] $M$ is Artinian.
    \item[(2)] $M$ is Noetherian.
    \item[(3)] $M$ is isomorphic to a finite direct sum of simple modules.
  \end{enumerate}
\end{exercise}

\begin{remark}
  From now on we mainly use the term ``semisimple module'' and not the term
  ``complemented module'',
  since they are the same thing.
\end{remark}

\subsection{Lattices*}
\label{ssec:lattices}
This subsection is outside the scope of the course,
but since we have use the modular property of the lattice of submodules
in \Cref{prop:properties-of-complemented-modules},
we might as well learn about it.

\begin{definition}[Lattice, as a poset]
  \label{def:lattice-poset}
  A partially order set $(L,\le)$ is called a lattice if any pair of elements $\set{a,b}$
  has a unique join (least upper bound) and a unique meet (greatest lower bound).
\end{definition}

\begin{remark}
  The join of $a$ and $b$ is denoted by $a \join b$ and $a \meet b$.
  We will soon show an example of a lattice,
  think why both the names, and the notations make sense (think what are
  the meet and the join in the lattice of subsets of $\R$ with respect to set
  inclusion).
\end{remark}

\begin{example}
  Consider the set of natural numbers where the order $a \le b$ is defined by
  $a \mid b$ ($a$ divides $b$).
  In this case, the join would be $a \join b = \lcm(a,b)$
  and the meet would be $a \meet v = \gcd(a,b)$.
\end{example}

\begin{example}[Lattice of subgroups]
  The subgroups of a group $G$ ordered by inclusion form a lattice,
  in which case we have $H \meet K = H \cap K$
  and $H \join K = \langle H \cup K \rangle$, for any two subgroups $H,K \leqslant G$.
\end{example}

\begin{example}[Lattice of submodules]
  The submodules of a module $M$ ordered by inclusion form a lattice,
  in which case we have $N \meet K = N \cap K$
  and $N \join K = N + K$, for any two submodules $N$, $K$ of $M$.
\end{example}

\begin{example}[Lattice of ideals]
  \label{ex:lattice-of-ideals}
  Let $\mathcal I(R)$ be the set of all ideals of the ring $R$.
  Under the partial order of set inclusion the bottom element
  (denoted by $\bot$) is the zero ideal $\set{0}$,
  and the top elements (denoted by $\top$) is the ring itself $R$.
  The join of ideals is $I \join J = I + J$,
  and the meet of ideals is $I \meet J = I \cap J$.
\end{example}

\begin{remark}
  We can also view lattices as algebraic structures.
\end{remark}

\begin{definition}[Lattice]
  \label{def:lattice-algebra}
  A lattice is an algebraic structure $(L,\join,\meet)$, consisting of a set $L$,
  and two binary, commutative and associative operations $\join$ and $\meet$ on $L$
  which satisfy the following axioms (sometimes called absorption laws):
  \[
    a \join (a \meet b) = a \qquad
    a \meet (a \join b) = a
  \]
\end{definition}

\begin{remark}
  From the axioms of the lattice follow the two ``idempotent laws'':
  \[
    a \join a = a \qquad a \meet a = a
  \]
  It also follows that both $(L,\join)$ and $(L,\meet)$ are semilattices
  (semilattices are defined by the structure induced on these structures by
  the bigger structure $(L,\join,\meet)$).
\end{remark}

\begin{exercise}
  \label{ex:think}
  Think why \Cref{def:lattice-poset} and \Cref{def:lattice-algebra} are equivalent
  definitions.
\end{exercise}

\begin{remark}[Bounded lattice]
  A bounded lattice is a lattice which has a top element, denoted by $1$ or $\top$,
  and a bottom element, denoted by $0$ or $\bot$.
  This makes $\mathcal I(R)$ (defined in \Cref{ex:lattice-of-ideals})
  a bounded lattice.
  Note that in bounded lattices the $0$ element is the identity element of the
  join operation, and the $1$ element is the identity for the meet operation.
  It can be shown that a partially ordered set is a bounded lattice if and only if
  every finite set of elements (including the empty set) has a join and a meet.
\end{remark}

\begin{remark}[Complete lattice]
  A poset is called a complete lattice if all its subsets have both a join and a meet.
  Hence every complete lattice is a bounded lattice.
\end{remark}

\begin{definition}[Distributive lattice]
  \label{def:distributive-lattice}
  A lattice $(L,\meet,\join)$ is said to be a distributive lattice
  if it satifies the following distributive axioms:
  \[
    a \join (b \meet c) = (a \join b) \meet (a \join c) \qquad
    a \meet (b \join c) = (a \meet b) \join (a \meet c)
  \]
\end{definition}

\begin{remark}
  It turns out that each distributive law in \Cref{def:distributive-lattice}
  implies the other.
\end{remark}

\begin{definition}[Modular lattice]
  A modular lattice is a lattice $(L,\meet,\join)$ which satisfies the
  modular law:
  \[
    a \le b \implies a \join (x \meet b) = (a \join x) \meet b,
  \]
  where $x,a,b$ are aribtrary elements in the lattice, $\le$
  is the partial order (see \Cref{ex:think}).
\end{definition}

\begin{remark}
  It turns out that the lattice of submodules of a module $M$ with respect
  to set inclusion is a modular lattice that is not distributive.
  In this case the modular property becomes:
  \[
    L \subseteq N \implies
    L + (K \cap N) = (L + K) \cap N,
  \]
  where $L,N,K$ are clearly submodules of $M$.
  By commutativity we get that:
  \[
    L \subseteq N \implies
    L + (N \cap K) = N \cap (L + K),
  \]
  which is exactly what we used in \Cref{prop:properties-of-complemented-modules}.
\end{remark}

\newpage

\section{Wedderburn structure theorem}

\begin{definition}[The $\alpha$ map setting]
Let $n \geq 2$, let $R$ be a ring, and let $M_1,\dots,M_n$
be right $R$-modules.
Let $M = \oplus_{i=1}^{n} M_i$,
and let $\rho_i \colon M_i \to M$ be the inclusions,
and $\pi_i \colon M \to M_i$ be the projections.
Moreover,
denote $\alpha_{i,j} := \pi_i \circ \alpha \circ \rho_j \in \Hom_R(M_j,M_i)$.
For $\alpha \in \End_R(M)$ and $\overline m \in M$ denote:
\[
  \alpha(\overline m) :=
  (\alpha_{i,j})_{1 \le i,j \le n}
  \begin{pmatrix}
    m_1 \\ \vdots \\ m_n
  \end{pmatrix} :=
  \begin{pmatrix}
    \sum_{j=1}^{n} \alpha_{1,j}(m_j) \\
    \vdots \\
    \sum_{j=1}^{n} \alpha_{n,j}(m_j)
  \end{pmatrix}
\]
Note the analogy with ordinary matrix-vector multiplication,
where the entries $\alpha_{i,j}$ are maps applied to the components $m_j$.
For $\alpha,\beta \in \End_R(M)$ denote:
\[
  (\alpha \circ \beta)_{1 \le i,j \le n} :=
  \sum_{k=1}^{n} \alpha_{i,k} \circ \beta_{k,j},
\]
which is analogous to matrix multiplication.
Conversely, for every $1 \le i,j \le n$ let $\gamma_{i,j} \in \Hom_R(M_j,M_i)$.
We can then define a map $\alpha \colon M \to M$ by:
for every $\overline m = (m_1,\dots,m_n)^T \in M$,
\[
  \alpha(\overline m) :=
  \begin{pmatrix}
    \sum_{j=1}^{n} \gamma_{1,j}(m_j) \\
    \vdots \\
    \sum_{j=1}^{n} \gamma_{n,j}(m_j)
  \end{pmatrix}.
\]
The map $\alpha$ is an $R$-endomorphism of $M$ and for every $1 \le i,j \le n$
we have $\alpha_{i,j} = \gamma_{i,j}$.
Let $E$ be the set of $n \times n$ matrices $X$ such that for every
$1 \le i, j \le n$, we have:
\[
  X_{i,j} \in \Hom(M_j,M_i).
\]
Let $\psi \colon \End_R(M) \to E$ be the map defined by:
for every $\alpha \in \End_R(M)$, we have:
\[
  \psi(\alpha) := (\alpha_{i,j})_{1 \le i,j \le n}.
\]
The map $\psi$ is an isomorphism of rings, where $E$ is endowed with matrix-like
operations:
\begin{itemize}
  \item Addition is defined as follows: for every $1 \le i,j \le n$
    and all elements $\alpha = (\alpha_{i,j})_{1 \le i,j \le n}$
    and $\beta = (\beta_{i,j})_{1 \le i,j \le n}$ of $E$,
    then:
    \[
      (\alpha + \beta)_{i,j} := \alpha_{i,j} + \beta_{i,j}.
    \]
  \item Multiplication is defined as follows: for every $1 \le i,j \le n$
    and all elements $\alpha = (\alpha_{i,j})_{1 \le i,j \le n}$
    and $\beta = (\beta_{i,j})_{1 \le i,j \le n}$ of $E$,
    then:
    \[
      (\alpha \beta)_{i,j} :=
      \sum_{k=1}^{n} \alpha_{i,k} \circ \beta_{k,j}.
    \]
\end{itemize}
\end{definition}

\begin{corollary}
  Let $M$ be an Artinian semisimple right $R$-module.
  Then:
  \[
    \End_R(M) \cong
    M_{n_1}(D_1) \times \cdots \times M_{n_k}(D_k).
  \]
\end{corollary}
\begin{proof}
  Since $M$ is semisimple and Artinian:
  \[
    M \cong \bigoplus_{i=1}^{k} M_i^{n_i},
  \]
  with simple and pairwise nonisomorphic modules
  (if the sum were infinite we could show $M$ is not Artinian
  by considering less and less coordinates).
  By Schur's lemma, $D_i := \End_R(M_i)$ is a division algebra,
  and for every $1 \le i \neq j \le n$, we have $\Hom_R(M_j,M_i) = 0$.
  Using the matrix description, we get that:
  \[
    \End_R(M) \cong
    M_{n_1}(D_1) \times \cdots \times M_{n_k}(D_k).
  \]
\end{proof}

\begin{definition}[Simple ring]
  A nonzero ring $R$ is called simple if $R$ and $0$ are the only two-sided
  ideals.
\end{definition}

\begin{example}
  A division ring is a simple ring.
\end{example}

\begin{example}
  We claim that if $R$ is a simple ring then $M_n(R)$ is a simple ring.
  \begin{enumerate}
    \item[(1)] We have to show that for every nonzero $X \in M_n(R)$,
      the two-sided ideal generated by $X$ is equal to $M_n(R)$.
    \item[(2)] Let $X \in M_n(R)$ be a nonzero element and fix $1 \le k,j \le n$
      such that $x := X_{k,l} \neq 0$.
    \item[(3)] For every $1 \le i,j \le n$ and every $t \in R$,
      let $e_{i,j}(t)$ be the matrix with $t$ on the $(i,j)$th entry
      and $0$ elsewhere.
    \item[(4)] In order to finish the proof,
      it is enough to show that for every $1 \le i, j \le n$,
      we have that $e_{i,j}(1)$ belongs to the two sided ideal generated by $X$.
    \item[(5)]
      Since $R$ is a simple ring there exist $r_1,s_1,\dots,r_m,s_m \in R$
      such that:
      \[
        1 = r_1 x s_1 + \cdots + r_m x s_m \implies
        e_{i,j}(1) =
        e_{i,k}(r_1) X e_{l,j}(s_1) + \cdots +
        e_{i,k}(r_m) X e_{l,j}(s_m),
      \]
      and this completes the proof.
  \end{enumerate}
\end{example}

\begin{exercise}
  \label{ex:simplicity-in-infinite-dimensions}
  Let $D$ be a field and $V$ an $\aleph_0$-dimensional $D$-module.
  Prove that:
  \begin{enumerate}
    \item[(1)] The set
    $I := \set{\alpha \in \End_D(V) \mid \dim(\im \alpha) < \infty}$
    is a nonzero proper two sided ideal of the space $E := \End_D(V)$.
    \item[(2)] $R := E/I$ is a simple ring.
    \item[(3)] $R_R$ does not contain a simple $R$-module.
  \end{enumerate}
\end{exercise}
\begin{proof} \phantom{}
  \begin{enumerate}
    \item[(1)] Note that a modules over a field is a vector space.
      Let $\alpha \in \End_D(V)$ be an endomorphism which projects
      onto a $1$ dimensional subspace of $V$.
      Since $\dim \im \alpha = 1 < \infty$ we have $\alpha \in I$
      hence $I \neq \emptyset$.

      Let $\alpha \in I$ and $\beta,\gamma \in E$.
      We have that $\im(\alpha \circ \beta) \subseteq \im(\alpha)$
      hence $\dim(\alpha \circ \beta) < \infty$
      and so $\alpha \circ \beta \in I$.
      We also have that $\im(\gamma \circ \alpha) = \gamma(\im \alpha)$,
      and since $\dim(\im \alpha) < \infty$, and $\gamma$
      is a linear map, we get that $\dim(\gamma \circ \alpha) < \infty$
      and hence $\gamma \circ \alpha \in I$.
      
      Additionally, $I$ is additively closed since
      $\im(\alpha + \beta) \subseteq \im \alpha + \im \beta$.

      The ideal $I$ is proper since $\dim(\id_V) = \aleph_0$
      hence $\id_V \notin V$.
    \item[(2)] By the correspondence theorem we must show that any ideal
      $J \supsetneq I$ in $E$ must be equal to $E$.
      Let $\alpha \in J \setminus I$,
      then $\dim \im \alpha = \aleph_0$,
      hence there exists some subspace $U \subseteq V$ such that the restriction
      $\alpha \mid_U \colon U \to \im \alpha$ is an isomorphism,
      hence $\dim U = \aleph_0$.
      Note that by the above claims, we have that
      there exist isomorphisms $\phi \colon V \to U$
      and $\psi \colon \im \alpha \to V$.
      Note that we have that:
      \[
        (\psi \circ P) \circ \alpha \circ \phi = \id_V,
      \]
      where $P \colon V \to \im \alpha$ is the projection map.
      Since both $\psi \circ P$ and $\phi$ are elements in $E$,
      we get that $\id_V \in \langle \alpha \rangle \subseteq J$,
      hence $J = E$, which completes the proof.
    \item[(3)] A simple submodule of $R_R$ corresponds to a minimal
      right ideal of $R$.
      We know from $(2)$ that $R$ has no nontrivial two sided ideals,
      but since it tells us nothing about the existence of right ideals,
      we have to do the work.

      Let $0 \neq \alpha \in R = E/I$ meaning $\im \alpha = \aleph_0$.
      Let $\im \alpha = W_1 \oplus W_2$
      so that $\dim W_1 = \dim W_2 = \aleph_0$.
      Let $\beta$ be the projection of $\alpha$ onto $W_1$,
      so that $\dim \im \beta = \dim W_1 = \aleph_0$ hence $0 \neq \beta \in R$.
      The inclusion $\beta R \subseteq \alpha R$ is obvious,
      and the inclusion is strict since $\alpha \notin \beta R$
      (because by construction we would get $\dim W_2 < \aleph_0$).
      This implies that there is not minimal right ideal in $R$,
      hence $R_R$ has no simple submodules.
  \end{enumerate}
\end{proof}

\begin{exercise}
  Let $n \geq 1$ and $D$ be a division ring.
  View the elements of $D^n$ as rows of elements in $R$.
  Prove that:
  \begin{enumerate}
    \item[(1)] $D^n$ is a simple right $M_n(D)$-module.
    \item[(2)] $\End_{M_n(D)}(D^n) \cong D$.
    \item[(3)] $\dim_{\End_{M_n(D)}(D^n)}(D^n) = n$.
  \end{enumerate}
\end{exercise}
\begin{proof} \phantom{}
  \begin{enumerate}
    \item[(1)] Let $0 \neq v \in D^n$.
      We must show that $v \cdot M_n(D) = D^n$.
      Since $D$ is a division ring, the conclusion is trivial.
    \item[(2)] 
    \item[(3)] 
  \end{enumerate}
\end{proof}

\begin{lemma}
  \label{lemma:submodules-of-rR-for-simple-R}
  Let $R$ be a simple ring.
  If $R_R$ contains a simple right submodule $M$ (equivalently, left),
  then all the simple submodules of $R_R$ are isomorphic to $M$ and $R_R$
  (equivalently, $_R R$) is semisimple.
\end{lemma}
\begin{proof}
  For every $r \in R$ we know that $rM$ is a submodule of $R_R$.
  If $rM \neq 0$, then the map $x \mapsto rx$ is a nonzero $R$-homomorphism
  from the simple right $R$-module $M$ onto the right $R$-module $rM$
  so it must be an isomorphism.
  Since $\sum_{r \in R} rM$ is a nonzero two sided ideal of the simple ring $R$,
  it must be equal to $R$.
  Thus, $R_R$ is an inner direct sum of the simple submodules isomorphic to $M$.
  In particular, $R$ is semisimple and every sime submodule of $R_R$ is
  isomorphic to $M$.
\end{proof}

\begin{definition}[$\rho$-isomorphism]
  Let $\rho \colon R \to S$ be an isomorphism of rings.
  Let $M$ be a right $R$-module and $N$ a right $S$-module.
  \begin{enumerate}
    \item[(1)] A bijective map $\psi \colon M \to N$ is called a
      $\rho$-isomorphism if for every $r \in R$ and $m \in M$,
      we have that $\psi(mr) = \psi(m) \rho(r)$.
    \item[(2)] We say that $M$ and $N$ are $\rho$-isomorphic if there exists
      a $\rho$-isomorphism between them.
  \end{enumerate}
\end{definition}

\begin{remark}
  Let $\rho \colon R \to S$ be an isomorphism of rings,
  and $\psi \colon M \to N$ a $\rho$-isomorphism of right modules.
  The map $\overline \psi \colon \End_R(M) \to \End_S(N)$ defined by
  $\overline \psi(\alpha) := \psi \circ \alpha \circ \psi^{-1}$
  for every $\alpha \in \End_R(M)$,
  is an isomorphism of rings.
  Moreover, $\psi$ is a $\overline \psi$-isomorphism from the left
  $\End_R(M)$-module $M$ to the left $\End_S(N)$ module $N$.
  Indeed, for every $\alpha \in \End_R(M)$ and every $m \in M$,
  we have:
  \[
    \psi(\alpha(r)) =
    \psi \circ \alpha \circ \psi^{-1}(\psi(m)) =
    \overline(\alpha)(\psi(m)).
  \]
\end{remark}

\begin{lemma}
  \label{lem:matrices-over-division-rings-are-unique}
  Let $n_1,n_2 \geq 2$ and $D_1$ and $D_2$ be division algebras.
  If $M_{n_1}(D_1) \cong M_{n_2}(D_2)$ then $D_1 \cong D_2$
  and $n_1 = n_2$.
\end{lemma}
\begin{proof}
  Let $\rho \colon M_{n_1}(D_1) \to M_{n_2}(D_2)$ be an isomorphism.
  For every $1 \le r \le 2$, we have that $M_{n_r}$ is a simple ring % why?
  so:
  \[
    L_r := \set{X \in M_{n_r}(D_r) \mid (\forall 2 \le i \le n_r \stand
    \forall 1 \le j \le n_r) X_{i,j} = 0}
  \]
  is a minimal right ideal of $M_{n_r}(D_r)$.
  Since $M_{n_r}(D_r)$ is simple,
  every minimal right ideal of $M_{n_r}(D_r)$ is isomorphic to $L_r$
  as an $M_{n_r}(D_r)$-module.
  Since $\psi$ sends a minimal right ideal of $M_{n_1}(D_1)$
  to a minimal right ideal of $M_{n_2}(D_2)$,
  we get that $L_1$ and $L_2$ are $\rho$-isomorphic,
  % check if theres notation=^p
  and:
  \[
    D_1 \cong
    \End_{M_{n_1}(D_1)}(L_1) \cong
    \End_{M_{n_2}(D_2)}(L_2) \cong
    D_2.
  \]
  Let
  $\overline \psi \colon \End_{M_{n_1}(D_1)}(L_1) \to \End_{M_{n_2}(D_2)}(L_2)$
  be the isomorphism induced by $\psi$.
  Since $L_1$ and $L_2$ are also $\overline \psi$-isomorphic:
  \[
    n_1 =
    \dim_{\End_{M_{n_1}(D_1)}(L_1)}(L_1) =
    \dim_{\End_{M_{n_2}(D_2)}(L_2)}(L_2) =
    n_2,
  \]
  which completes the proof.
\end{proof}

\begin{example}
  \Cref{lem:matrices-over-division-rings-are-unique} is false over
  general rings.
  For example, $M_4(\mathbb F) \cong M_2(M_2(\mathbb F))$.
  %% explain.
\end{example}

\begin{lemma}
  \label{lem:unique-factorization-of-product-modules}
  Let $R_1,\dots,R_m$ and $S_1,\dots,S_n$ be simple rings and denote
  $R := R_1 \times \cdots \times R_n$ and $S := S_1 \times \cdots \times S_n$,
  and suppose that $R \cong S$.
  Then $m = n$ and there exists a permutation $\sigma \in \Sym(n)$
  such that for every $1 \le i \le m$,
  the map $\psi$ induces an isomorphism of rings
  $\overline \psi_i \colon R_i \to S_{\sigma(i)}$.
\end{lemma}
\begin{proof}
  We can view $R$ (respectively, $S$) as an inner direct product of the $R_i$s
  (respectively, of the $S_j$s).
  In particular, we identify each $R_i$ (respectively, each $S_j$)
  as a two sided ideal of $R$ (respectively, $S$).
  Since each $R_i$ (respectively, each $S_j$)
  is a simple ring,
  the sets $X := \set{R_i \mid 1 \le i \le m}$
  and $Y := \set{S_j \mid 1 \le j \le n}$ are the sets of minimal
  two sided ideals of $R$ and $S$ respectively.
  Moreover, $|X| = m$ and $|Y| = n$.
  Since $\psi$ is an isomorphism of rings, $m = n$ and there exists
  $\sigma \in \Sym(m)$ such that for every $1 \le i \le m$,
  we have $R_i = S_{\sigma(i)}$. % is this wrong?
  Let $1 \le i \le m$ and let $\overline \psi_i \colon R_i \to S_{\sigma(i)}$
  be the map induced by $\psi$.
  Since $\psi$ is a homomorphism of rings, $\overline \psi$ preserves
  addition and multiplication.
  It is left to show that for every $1 \le i \le m$,
  we have $\beta_i(1_{R_i}) = 1_{S_{\sigma(i)}}$.
  Since the $\beta_i$ respect multiplication and are surjective,
  $\beta_i(1_{R_i})$ is an element of the ring $S_{\sigma(i)}$ such that
  for every $s \in S_{\sigma(i)}$ we have $s \beta_i(1_{R_i}) = s$.
  In particular, by taking $s := 1_{S_{\sigma(i)}}$ we get that:
  \[
    \beta_i(1_R) =
    1_{S_{\sigma(i)}} \beta_i(1_{R_i}) =
    1_{S_{\sigma(i)}},
  \]
  which completes the proof.
\end{proof}

\begin{corollary}
  Let $m_1,\dots,m_r$ and $n_1,\dots,n_s$ be positive integers,
  and let $D_1,\dots,D_r$ and $E_1,\dots,E_s$ be division algebras
  such that:
  \[
    M_{m_1}(D_1) \times \cdots \times M_{m_r}(D_r) \cong
    M_{n_1}(E_1) \times \cdots \times M_{n_s}(E_s).
  \]
  Then $r = s$ and there exists a permutation $\sigma \in \Sym(r)$
  such that for every $1 \le i \le r$ we have $n_{\sigma(i)}$
  and $D_i \cong E_{\sigma(i)}$.
\end{corollary}
\begin{proof}
  Immediate by \Cref{lem:matrices-over-division-rings-are-unique}
  and \Cref{lem:unique-factorization-of-product-modules}.
\end{proof}

\begin{definition}
  A ring $R$ is called a right (equivalently, left)
  Artinian/Noetherian/semisimple ring if
  $R$ is Artinian/Noetherian/semisimple as a right (equivalently, left)
  $R$-module.
\end{definition}

\begin{example}
  Here are some claims based on the previous definition:
  \begin{enumerate}
    \item[(1)] The ring $\Z$ is not semisimple.
    \item[(2)] The ring $R := \Z/6\Z$ is not simple.
      It is semisimple since $R_R \cong \Z/2\Z \oplus \Z/3\Z$.
    \item[(3)] In \Cref{ex:simplicity-in-infinite-dimensions},
      we constructed a simple ring
      that does not contain a minimal ideal.
      It follows that this ring is not semisimple.
    \item[(4)] For every $n \geq 1$ and any division ring $D$,
      we have that $M_n(D)$ is right semisimple.
      Indeed, as a right module it is isomorphic to a direct sum of $n$
      copies of $D^n$ where there elements of $D^n$ are viewed as rows
      and $M_n(D)$ act by multiplication from the right.
    \item[(5)] Similarly, for every $n \geq 1$ and any division ring
      $D$ we have that $M_n(D)$ is right semisimple.
    \item[(6)] If $R_1,\dots,R_n$ are semisimple rings,
      so is $R_1 \times \cdots \times R_n$. % prove this! <- Chen comment
  \end{enumerate}
\end{example}

\begin{lemma} \phantom{}
  \begin{enumerate}
    \item[(1)] A right semisimple ring is right Artinian.
    \item[(2)] A left semisimple ring is left Artinian.
    \item[(3)] A simple ring is right (equivalently, left) semisimple
      if and only if it is right (equivalently, left) Artinian.
  \end{enumerate}
\end{lemma}
\begin{proof} \phantom{}
  \begin{enumerate}
    \item[(1)] The module $R_R$ is generated by its simple submodules.
      In particular, there are finitely many simple submodules
      $M_1,\dots,M_k$ such that $1_R \in \sum_{i=1}^{k} M_k$.
      It follows that
      $R = \set{1_R r \mid r \in R} \subseteq \sum_{i=1}^{k} M_k$.
      Since $R_R$ is generated by finitely many simple submodules,
      $R$ is right Artinian.
    \item[(2)] Analogous to $(1)$.
    \item[(3)] By $(1)$, if $R$ is semisimple it is Artinian.
      Assume that $R$ is simple and Artinian.
      Since $R$ is Artinian, $R_R$ contains a minimal submodule $M$.
      This claim follows from \Cref{lemma:submodules-of-rR-for-simple-R}.
  \end{enumerate}
\end{proof}

\begin{theorem}[Wedderburn structure theorem]
  Let $R$ be a ring.
  The following properties are equivalent:
  \begin{enumerate}
    \item[(1)] $R$ is a right semisimple ring.
    \item[(2)] $R$ is a left semisimple ring.
    \item[(3)] There exist positive integers $n_1,\dots,n_k$
      and division algebras $D_1,\dots,D_k$ such that
      $R \cong M_{n_1}(D_1) \times \cdots \times M_{n_k}(D_k)$.
  \end{enumerate}
  Moreover, if $(3)$ holds, then $n_1,\dots,n_k$ are unique and
  $D_1,\dots,D_k$ are unique up to isomorphism.
\end{theorem}
\begin{proof}
  
\end{proof}

\begin{corollary}
  An Artinian simple ring is isomorphic to $M_n(D)$
  for unique $n \geq 2$ and unique (up to isomorphism) division algebra $D$.
\end{corollary}

\begin{corollary}
  There exists a simple ring which is not semisimple.
\end{corollary}
\begin{proof}
  The ring $R = E/I$ defined in \Cref{ex:simplicity-in-infinite-dimensions}
  is simple (as shown in part $(2)$ of the exercise),
  but since it is shown in part $(3)$ that it is not Artinian,
  it cannot satisfy condition $(3)$ in the Wedderburn structure theorem,
  hence it cannot be semisimple.
\end{proof}

\newpage

\section{Radicals}

\begin{definition}[Radical of a module]
  Let $R$ be a ring and $M$ be a left (or right) $R$-module.
  The radical of $M$ is defined to be the intersection of all maximal
  modules of $M$.
  If none exist, then we define it to be $M$.
\end{definition}

\begin{remark}
  We denote the radical of $M$ by $\Rad M$.
\end{remark}

\begin{example}
  We have the following radicals:
  \begin{enumerate}
    \item[(1)] $\Rad \Q_{\Q} = 0$.
    \item[(2)] $\Rad \Q_{\Z} = \Q$.
    \item[(3)] $\Rad(\Z/6\Z) = 2\Z/6\Z \cap 3\Z/6\Z = 0$.
    \item[(4)] $\Rad \Z/18\Z = 2\Z/18\Z \cap 3\Z/18\Z = 6\Z/18\Z$.
  \end{enumerate}
\end{example}

\begin{exercise}
  Let $n \geq 1$, and let $D$ be a division ring.
  Denote by $R = M_n(D)$.
  Show that $\Rad _R R = \Rad R_R = 0$.
\end{exercise}

\begin{lemma}
  \label{lem:radicals}
  Let $R$ be a ring and $\alpha \colon M \to N$ an $R$-homomorphism.
  Then $\alpha(\Rad M) \subseteq \Rad N$.
\end{lemma}
\begin{proof}
  We have that:
  \[
    \alpha^{-1}(\Rad N) =
    \bigcap_{L \text{ maximal in } N} \alpha^{-1} L,
  \]
  and thus we simply need to show that if $L$ is a maximal module of $N$,
  then $\alpha^{-1}(L)$ contains $\Rad M$.
  Let $L$ be a maximal submodule of $N$.
  $\alpha$ induces an injective $R$-homomorphism:
  \[
    \bar \alpha \colon M/\alpha^{-1}(L) \to N/L,
  \]
  where $N/L$ is a simple ideal,
  hence there are only two possiblities:
  \begin{enumerate}
    \item[(1)] $\im \bar \alpha = 0$ and then $\alpha^{-1}(L) = M$.
    \item[(2)] $\bar \alpha$ is an isomorphism, and hence
      $P := \alpha^{-1}(L)$ is a maximal submodule of $M$
      and so $\Rad M \subseteq P = \alpha^{-1}(L)$.
  \end{enumerate}
  This completes the proof.
\end{proof}

\begin{corollary}
  Let $R$ be a ring.
  Then $\Rad R_R$ and $\Rad _R R$ are two-sided ideals.
\end{corollary}
\begin{proof}
  Without loss of generality we will only prove the case for $\Rad R_R$.
  It is clear that $\Rad R_R$ is a right ideal.
  Thus it suffices to show that for every $r \in R$,
  \[
    r \Rad R_R \subseteq \Rad R_R.
  \]
  This follows from \Cref{lem:radicals}
  since the map $R_R \to R_R$ which sends $x \mapsto r x$
  is an $R$-homomorphism.
\end{proof}

\begin{lemma}
  Let $R$ be a ring and $M$ be an $R$-module.
  Then for every $m \in M$ the following conditions are equivalent:
  \begin{enumerate}
    \item[(1)] $m \in \Rad M$.
    \item[(2)] For every $N \leqslant M$,
      if $M$ is generated by $N \cup \set{m}$ then $N = M$.
  \end{enumerate}
\end{lemma}
\begin{proof} \phantom{}
  \begin{enumerate}
    \item[$(1) \implies (2)$] Let $N$ be a submodule of $M$
      such that $N \cup \set{m}$ generates $M$.
      Assume, by way of contradiction, that $M = N$.
      Since $N \cup \set{m}$ generates $M$,
      we conclude that $m \notin N$, hence $M \neq N$.
      Denote:
      \[
        X := \set{L \leqslant M \mid N \subseteq L \stand m \notin L}.
      \]
      The set $X$ is nonempty since $N \in L$ and it is closed under unions
      of increasing chains, and hence by Zorn's lemma there exists a maximal
      element $P$.
      We know that $M \notin X$ and hence $P \neq M$.
      $P$ is a maximal submodule with respect to the property of containing $N$
      and not containing $m$.
      We will now show that $P$ is a maximal submodule of $M$
      in contradiction to:
      \[
        P \supseteq \Rad M \stand m \notin P \stand m \notin \Rad M,
      \]
      where the first inclusion follows from $P$ being a maximal submodule.
      Assume that exists $P \leqslant Q \leqslant M$ with $P \neq Q$.
      Then since $P$ is maximal in $X$, we get that $Q \notin X$,
      but $N \subseteq P \subseteq Q$ and so $m \in Q$.
      This implies that $Q = M$ since $M$ is generated by $N \set{m}$.
    \item[$(2) \implies (1)$]
      Let $m \in M$ be an element that satisfies the property in $(2)$.
      Let $L$ be a maximal submodule of $M$.
      We need to show that $m \in L$.
      Assume, by way of contradiction, that $m \notin L$.
      Then, by maximality of $L$,
      the submodule generated by $L \cup \set{m}$ is $M$.
      By assumption on $m$, we get $M = L$.
      This contradiction completes the proof.
  \end{enumerate}
\end{proof}

\begin{lemma}[Nakayama's lemma]
  Let $R$ be a ring, let $M$ be an $R$-module,
  and let $P \leqslant M$.
  Then:
  \begin{enumerate}
    \item[(1)] Assume that for every $N \leqslant M$,
      if $P + N = M$ then $N = M$.
      Then $P \subseteq \Rad M$.
    \item[(2)] If $P \subseteq \Rad M$ and $M$ or $P$
      are finitely generated as modules over $R$,
      then for every submodule $N \leqslant M$,
      if $M = P + N$ then $M = N$.
  \end{enumerate}
\end{lemma}
\begin{proof} \phantom{}
  \begin{enumerate}
    \item[(1)] We need to show that if $N$ is a maximal submodule of $M$,
      then $P \leqslant N$.
      Since $N \neq M$ (by definition) the assumption of $(1)$
      shows that $P + N \neq M$.
      From maximality of $N$, we get $N = P + N$,
      which means that $P \leqslant N$.
    \item[(2)] If $P$ is finitely generated as a module over $R$,
      then we get the claim from induction over the minimal number of
      generators of $P$ and the previous lemma.
      Next assume that $M$ is finitely generated and that $N \leqslant M$
      such that $M = P + N$.
      Denote the generators of $M$ by $x_1,\dots,x_k$.
      Since $M = P + N$ we get that for any $1 \le i \le k$ there exists
      $y_i \in P$ and $z_i \in N$ such that $x_i = y_i + z_i$.
      Denote by $Q$ the submodules of $M$ generated by $y_1,\dots,y_k$.
      Then:
      \begin{enumerate}
        \item[(1)] $Q \subseteq P \subseteq \Rad M$.
        \item[(2)] $Q$ is finitely generated.
        \item[(3)] $M = Q + N$.
      \end{enumerate}
      From the previous case we get that $M = N$, which completes the proof.
  \end{enumerate}
\end{proof}

\begin{example}
  We have that $P = \Rad \Q_{\Z} = \Q$.
  We have $0 \neq \Q$ and $0 + P = \Q$.
  This shows that the assumption of being finitely generated is necessary.
\end{example}

\subsection{Jabobson's radical}

\begin{lemma}[Nakayama's lemma]
  \label{lem:nakayama}
  Let $R \neq 0$ be a ring with maximal $P \idealin R$.
  Then the following conditions are equivalent:
  \begin{enumerate}
    \item[(1)] $P \subseteq \Rad R_R$.
    \item[(2)] If $M$ is a finitely generated right $R$-module
      and $N \leqslant M$ satisfies $M = N + MP$,
      then $N = M$.
    \item[(3)] $\set{1 + x \mid x \in P} \in R^{\times}$.
  \end{enumerate}
\end{lemma}
\begin{proof} \phantom{}
  \begin{enumerate}
    \item[$(1) \implies (2)$]
      Let $M$ be a finitely generated module over $R$.
      For every $m \in M$ the map $\alpha_m \colon R_R \to M$
      given by $x \mapsto mx$ is an $R$-homomorphism.
      Thus, by the lemma from the beginning of the section:
      \[
        MP =
        \sum_{m \in M} mP =
        \sum_{m \in M} \alpha_m(P) \subseteq
        \Rad M.
      \]
      The rest follows from Nakayama's lemma for modules.
    \item[$(2) \implies (3)$] Denote $G := \set{1 + x \mid x \in P}$.
      Since $P$ is a right ideal,
      $G$ is closed under multiplication and contains $1$.
      We will show that every element in $G$ has a right inverse in $R$.
      First, Assume by way of contradiction, that $x \in P$
      and that $1 + x$ does not have a right inverse in $R$.
      Then $I := (1+x) R \neq R$ is a right ideal and hence contained in a
      maximal right ideal $J$.
      Since $1 + x \in J$ then $x \notin J$.
      Hence, since $J$ is maximal $R_R = J + xR$.
      From $(2)$ it follows that $J = R_R$,
      which is a contradiction.

      We finished showing that for any $g \in G$ there exists a right inverse
      in $R$.
      Now we will prove that any $g \in G$ has a right inverse in $G$.
      Let $g \in G$.
      There exist $x \in P$ and $y \in R$ such that $g = 1 + x$
      and $1 = xy$.
      Denote $z := y - 1$, and then:
      \[
        1 = (1 + x) (1 + z) = 1 + x + z + xz,
      \]
      and $z = -x - xz \in P$.
      Thus, $y = 1 + z \in G$.
      We will now show that for any $g \in G$ there is a left inverse in $G$
      which completes this part of the proof.
      Let $g \in G$, there exists $h \in G$ inverse to $g$
      % TO BE CONTINUED
    \item[$(3)$]
  \end{enumerate}
\end{proof}

\begin{corollary}
  Let $R$ be a ring.
  Then $\Rad _R R = \Rad R_R$.
\end{corollary}
\begin{proof}
  By the left version of Nakayama's lemma for rings (specifically $(3)$
  implies $(1)$),
  we get that $\set{1 + x \mid x \in \Rad_R R} \subseteq R^{\times}$.
  We have seen before that $\Rad _R R$ is a two sided ideal,
  and in particular a right ideal.
  By $(3)$ implies $(1)$ (in the right versino of Nakayama's lemma)
  we get that $\Rad _R R \subseteq \Rad R_R$.
  From symmetry, we get that other direction of containment,
  hence they are equal.
\end{proof}

\begin{definition}[Jacobson radical]
  Let $R$ be a ring.
  The Jacobson radical of $R$ is defined by:
  \[
    J(R) := \Rad _R R = \Rad R_R.
  \]
\end{definition}

\begin{exercise}
  Show that if $R$ is an Artinian ring, then $R/J(R)$
  is a semisimple ring.
\end{exercise}

\begin{exercise}
  Let $R$ be a ring.
  Let $I$ be a right (or left ideal),
  such that all of its elements are nilpotent.
  Show that $I \in J(R)$.
\end{exercise}

\subsection{Radicals in Artinian rings}

\begin{proposition}
  Let $R$ be an Artinian ring.
  Then there exists $k \geq 1$ such that $J(R)^k = 0$.
\end{proposition}
\begin{proof}
  We know that:
  \[
    \cdots \subseteq J(R)^3 \subseteq J(R)^2 \subseteq J(R).
  \]
  Since $R$ is Artinian, there exists $k \geq 1$
  such that $J(R)^k = J(R)^{k+m}$ for every $m \in \N$.
  In particular, $J(R)^k = J(R)^{2k}$.
  We claim that $J(R)^k = 0$.
  Denote:
  \[
    X := \set{0 \neq I \idealin R_R \mid I J(R)^k = I}.
  \]
  The set $X$ is nonempty since $J(R)^k \in X$.
  Since $X$ is Artinian, there exists a minimal element $L$ in $X$.
  Since $L \neq 0$, there exists $0 \neq \ell \in L$.
  Note that:
  \[
    L \supseteq \ell J(R)^k \in X,
  \]
  so by minimality of $L$,
  \[
    L = \ell J(R)^k \subseteq \ell R \subseteq L,
  \]
  and so $L$ is finitely generated as an ideal.
  Moreover, $L = 0 + L = 0 + L J(R)^k$.
  By Nakayama's lemma for rings we get that $L = 0$,
  which is a contradiction.
\end{proof}


\end{document}
